%% kapa: exam=1 round=23 year=2012 period=2 form=A subject=appraisal_law subject_ko=부동산관계법규 seq=14 number=54
%% source: Q-Net 감정평가사 기출문제 (cst00309) artlSeq=5052580 file=raw/1008/5052580/2012년 제23회 감정평가사 제1차시험 기출문제(2교시 A형).pdf
%% conversion: kordoc 4.24.51 -> tools/gen1.py (auto); review=proofread-sonnet; flags=none
\begin{question}{54}
국토의 계획 및 이용에 관한 법령상 지형도면에 관한 설명으로 옳은 것은?
\begin{choices}
  \item 지형도면은 도시관리계획에 대한 주민의견청취 이전에 작성ㆍ제공되어야 한다.
  \item 지적이 표시된 지형도의 데이터베이스가 구축되어 있는 경우에도 이를 사용하여 지형도면을 작성할 수 없다.
  \item 도시관리계획 결정의 효력은 원칙적으로 지형도면에 대한 도지사의 승인ㆍ고시로 발생한다.
  \item 도시지역 외의 지역으로서 해당 토지에 도시계획시설이 결정되지 아니한 경우에는 지적이 표시되지 아니한 축척 5천분의 1 이상의 지형도에 도시관리계획사항을 명시한 도면을 작성할 수 있다.
  \item 주거지역에서 축척 3천분의 1 이상의 지형도를 사용하여 도시관리계획결정을 고시한 경우에는 지형도면을 따로 작성하지 않을 수 있다.
\end{choices}
\end{question}
\answer{4}
